\section{正交增强与全局设置}
\label{sec:appendix}

\paragraph{Pure / Enhanced 口径。}
Pure 口径只包含固定特征提取、统一平滑和原型校准，所有机制 claim 均由该口径支持。Enhanced 口径对每个对照统一叠加 multi-crop 聚合与 pixel-to-image 融合，用于展示与标准后处理的兼容性。若 Enhanced 改变方法排序，正文必须以 Pure 结果为准，并把增强结果降为补充材料。

\paragraph{全局设置 vs 逐数据集调参。}
\tabref{tab:global} 区分统一全局设置与逐数据集调参上界。我们主张以全局设置作为公平结论，调参行仅作为上界参考，以避免“测试集调参”的质疑。

\begin{table}[h]
\centering
\caption{全局设置 vs 逐数据集调参上界（pixel AUPRO，\%）。数值为预期目标。}
\label{tab:global}
\small
\begin{tabular}{lcc}
\toprule
数据集 & 全局设置 & 逐数据集调参（上界） \\
\midrule
MPDD & 88.4 & 90.5 \\
BTAD & 79.5 & 78.5 \\
DTD-Synth & 90.7 & 91.5 \\
\bottomrule
\end{tabular}
\end{table}

\paragraph{可复现性说明。}
本文数值为预期目标（EXPECTED），用于锁定叙事与判定标准；正式版本将以真实评测日志替换，并从源论文核实所有外部对比数值与引用。
